\section{CAT3D：多视角扩散模型 + NeRF 蒸馏}

\subsection{方法概述}

CAT3D 是一种可以接受文本提示、单张图像或多张图像作为输入，生成完整3D场景的方法。其核心思想是将多视角扩散模型与 NeRF 重建流程相结合。

\subsection{模型架构}

CAT3D 的输入包含两类视角：

\begin{itemize}
    \item \textbf{已观察视角（Observed Views）}：包含输入图像及其对应的相机参数。图像通过预训练的 VAE 编码为图像隐变量（image latent），相机信息编码为逐像素的射线坐标图（per-pixel ray coordinate map），包含每个像素的射线原点和方向。
    \item \textbf{新视角（Novel Views）}：没有真实图像，取而代之的是噪声隐变量，同时附带目标视角的相机位姿信息。
\end{itemize}

为帮助模型区分两类视角，还会输入一个\textbf{掩码}（mask），标识哪些视角是已观察的、哪些是待生成的。

\begin{knowledgebox}{射线坐标图（Ray Coordinate Map）}
射线坐标图是一种将相机位姿信息编码为空间表示的方式。对于图像中的每个像素，计算其对应的射线原点（ray origin）和射线方向（ray direction），将这些信息存储为与图像尺寸相同的坐标图。这种表示方式允许模型直接在像素级别理解3D几何关系。
\end{knowledgebox}

模型在所有帧之间使用\textbf{注意力机制}（attention），因此能够生成多视角一致的输出。训练数据为大量带位姿标注的多视角图像，监督模型在给定已知视角的条件下生成一致的新视角。

\subsection{两阶段流程}

CAT3D 采用两阶段流程：

\begin{enumerate}
    \item \textbf{多视角生成}：给定少量输入视角，使用多视角扩散模型生成场景的\textbf{密集}视角集
    \item \textbf{NeRF 重建}：将生成的所有视角输入标准 NeRF 管线进行3D重建
\end{enumerate}

\begin{warningbox}{生成的多视角并非完美3D一致}
多视角扩散模型生成的图像并非完全3D一致——不同视角之间可能存在微小的不一致。但关键洞察是：这些图像\textbf{足够一致}，使得 NeRF 重建管线能够将它们整合为一个一致的3D输出。NeRF 本质上充当了一个「平均化」的角色，将多视角中的不一致性消解掉。
\end{warningbox}

\subsection{性能}

端到端地，CAT3D 能够在约 \textbf{1分钟} 内从任意数量的输入（甚至仅一张图像）生成完整的3D场景。

\subsection{本章小结}

CAT3D 验证了「2D生成 + 3D蒸馏」范式的可行性：多视角扩散模型负责高质量图像生成，NeRF 负责将这些图像蒸馏为一致的3D表示。但 NeRF 优化步骤仍然是推理速度的瓶颈。

